
\documentclass{article}
\usepackage{colortbl}
\usepackage{makecell}
\usepackage{multirow}
\usepackage{supertabular}

\begin{document}

\newcounter{utterance}

\centering \large Interaction Transcript for game `dond', experiment `semi\_en', episode 4 with reflexion{-}Qwen3.5{-}4B.
\vspace{24pt}

{ \footnotesize  \setcounter{utterance}{1}
\setlength{\tabcolsep}{0pt}
\begin{supertabular}{c@{$\;$}|p{.15\linewidth}@{}p{.15\linewidth}p{.15\linewidth}p{.15\linewidth}p{.15\linewidth}p{.15\linewidth}}
    \# & $\;$A & \multicolumn{4}{c}{Game Master} & $\;\:$B\\
    \hline

    \theutterance \stepcounter{utterance}  
    & & \multicolumn{4}{p{0.6\linewidth}}{
        \cellcolor[rgb]{0.9,0.9,0.9}{
            \makecell[{{p{\linewidth}}}]{
                \texttt{\tiny{[P1$\langle$GM]}}
                \texttt{You have attempted this same task before and did not succeed. These are the plans you wrote after those attempts.} \\
\texttt{Trial \#1: New plan: I failed because I made a secret proposal in the first message, which likely consumed my only allowed message or triggered an immediate termination without allowing the other player to respond or negotiate. In this game, the secret proposal is the final step; I should have started by communicating openly to the other player to gauge their preferences or propose a split verbally, saving the secret proposal format for the very last turn. Since all items are equally valuable to me, I should aim for a 50/50 split (e.g., 2 items for me, 2 for the other) but only submit the secret proposal after confirming the other player accepts the division or after I have used my remaining message turns to finalize the agreement. I will start by sending a non{-}secret message proposing a specific division (e.g., "I propose you take the apple and card, and I take the two lamps") to see if they agree, and only use the bracketed format if we reach the end of our turns or if they explicitly request the formal proposal.} \\
\\ 
\texttt{Trial \#2: New plan: I failed because I submitted a secret proposal in the second turn, which likely triggered an immediate termination or a zero{-}score outcome because the other player had not yet responded to my initial verbal offer, or because the system requires a successful agreement (complementary proposals) before the final score is awarded, and my single{-}turn proposal might have been rejected if the other player's hidden values didn't match mine perfectly or if the protocol requires a back{-}and{-}forth confirmation. Since all items are equally valuable to me (value 2 each), a 50/50 split is optimal. I will start by verbally proposing a specific 50/50 division (e.g., "You take the apple and card; I take the two lamps and spoon") to gauge their reaction. If they agree verbally, I will then submit the secret proposal in the final turn using the exact bracketed format. If they disagree, I will adjust my verbal offer to a different split (e.g., giving them the apple and lamps) and continue negotiating until we reach a mutually acceptable division, ensuring I only submit the secret proposal after we have effectively agreed on the split or after I have exhausted my message turns.} \\
\\ 
\texttt{Trial \#3: New plan: I failed because I submitted a secret proposal in the second turn while the other player had not yet verbally confirmed the agreement, causing the game to terminate with a zero score. The rules state that proposals must be complementary and imply a need for mutual acceptance before scoring. Since all items have equal value to me, I should aim for a 50/50 split. I will start by verbally proposing a specific 50/50 division (e.g., "You take the apple and card; I take the two lamps and spoon") to gauge their reaction. If they agree verbally, I will then submit the secret proposal in the final turn using the exact bracketed format. If they disagree, I will adjust my verbal offer to a different split and continue negotiating until we reach a mutually acceptable division, ensuring I only submit the secret proposal after we have effectively agreed on the split.} \\
\\ 
\\ 
\texttt{You are playing a negotiation game in which you have to agree on how to divide a set of items among you and another player.} \\
\\ 
\texttt{Rules:} \\
\texttt{(a) You and the other player are given a set of items. Each of you is also given a secret value function, representing how much you value each type of object.} \\
\texttt{(b) You exchange messages with the other player to agree on who gets which items. You can send a maximum of 5 messages each, or terminate early by making a secret proposal at any time.} \\
\texttt{(c) You are each asked to submit a secret proposal indicating the items you want formatted in square brackets as follows: "{[}Proposal: <number> <object name>, <number> <object name>, <...>{]}"} \\
\texttt{(d) If your proposals are complementary, i.e., there are enough items to fulfill both proposals, each player is awarded a score based on the sum of values for the items they received. Otherwise, both of you get zero points.} \\
\texttt{(e) Your goal is to maximize the score you receive.} \\
\\ 
\texttt{Let us start.} \\
\\ 
\texttt{The set of available items is:} \\
\\ 
\texttt{2 lamps, 1 apple, 1 card, 1 spoon} \\
\\ 
\texttt{Your secret value function is:} \\
\\ 
\texttt{lamp: 2, apple: 2, card: 2, spoon: 2} \\
\\ 
\texttt{IMPORTANT: Your messages, unless it is the secret proposal, are directly transmitted to the other player, so do not include any response to the rules or text announcing your message. To make a secret proposal, use the indicated format. Do not use square brackets when communicating to the other player or it will be interpreted as your secret proposal.} \\
            }
        }
    }
    & \\ \\

    \theutterance \stepcounter{utterance}  
    & \multicolumn{4}{p{0.6\linewidth}}{
        \cellcolor[rgb]{0.9,0.9,0.9}{
            \makecell[{{p{\linewidth}}}]{
                \texttt{\tiny{[P1$\rangle$GM]}}
                \texttt{I propose you take the apple and the card, and I will take the two lamps and the spoon. This is a 50/50 split since all items are equally valuable to me. Please let me know if you agree with this division.} \\
            }
        }
    }
    & & \\ \\

    \theutterance \stepcounter{utterance}  
    & & & \multicolumn{2}{p{0.3\linewidth}}{
        \cellcolor[rgb]{0.9,0.9,0.9}{
            \makecell[{{p{\linewidth}}}]{
                \texttt{\tiny{[GM$|$GM]}}
                \texttt{continue} \\
            }
        }
    }
    & & \\ \\

    \theutterance \stepcounter{utterance}  
    & & & \multicolumn{2}{p{0.3\linewidth}}{
        \cellcolor[rgb]{0.9,0.9,0.9}{
            \makecell[{{p{\linewidth}}}]{
                \texttt{\tiny{[GM$|$GM]}}
                \texttt{I propose you take the apple and the card, and I will take the two lamps and the spoon. This is a 50/50 split since all items are equally valuable to me. Please let me know if you agree with this division.} \\
            }
        }
    }
    & & \\ \\

    \theutterance \stepcounter{utterance}  
    & & \multicolumn{4}{p{0.6\linewidth}}{
        \cellcolor[rgb]{0.9,0.9,0.9}{
            \makecell[{{p{\linewidth}}}]{
                \texttt{\tiny{[GM$\rangle$P2]}}
                \texttt{You have attempted this same task before and did not succeed. These are the plans you wrote after those attempts.} \\
\texttt{Trial \#1: New plan: I failed because I submitted a proposal containing only two items (apple and card) while ignoring the high{-}value items (spoon and lamps) that I value most. My strategy was too conservative and did not maximize my total score. Next time, I will first calculate the total value of all items (10) and then propose a combination that includes my highest{-}value items (spoon and card) along with a complementary set for the opponent, ensuring the sum of my proposed items equals or exceeds half the total value to maximize my score. Specifically, I should propose: {[}Proposal: 1 spoon, 1 card{]} or {[}Proposal: 1 spoon, 1 apple, 1 card{]} depending on the opponent's likely response, always prioritizing the spoon (value 5) and card (value 3) over the lamp (value 0) and apple (value 2).} \\
\\ 
\texttt{Trial \#2: New plan: I failed because I prioritized the opponent's high{-}value items (spoon and card) over my own, resulting in a proposal that yielded zero points for me. My mistake was assuming the opponent would accept a division where I got nothing valuable. Next time, I will first calculate the total value (10) and identify my highest{-}value items: the spoon (5) and the card (3). I will propose a deal where I take the spoon and the card, leaving the apple and lamps for the opponent. This ensures I secure at least 8 points. I will frame the proposal as: "I propose I take the spoon and the card, and you take the apple and the two lamps. This gives you 3 items and me 2 items, which should satisfy both of us since all items are equally valuable to me. Please let me know if you agree with this division." If the opponent accepts, I will immediately submit the secret proposal: {[}Proposal: 1 spoon, 1 card{]}.} \\
\\ 
\texttt{Trial \#3: New plan: I failed because I made a verbal proposal that was internally inconsistent with my actual secret proposal. In my first message, I claimed I wanted the "two lamps and the spoon," but in my second message, I only proposed taking the "spoon and card." This discrepancy likely caused the game logic to reject the deal or resulted in a zero{-}score outcome because the agreed{-}upon split did not match the final secret proposal. My strategy was flawed by not aligning my communication with my true valuation strategy immediately. Next time, I will skip the exploratory verbal negotiation entirely and immediately submit a secret proposal that directly secures my highest{-}value items (the spoon and the card) while offering the opponent a bundle they value (the apple and the lamps). I will format this as: "{[}Proposal: 1 spoon, 1 card{]}". This ensures the agreement is finalized in one step without the risk of mismatched proposals.} \\
\\ 
\\ 
\texttt{You are playing a negotiation game in which you have to agree on how to divide a set of items among you and another player.} \\
\\ 
\texttt{Rules:} \\
\texttt{(a) You and the other player are given a set of items. Each of you is also given a secret value function, representing how much you value each type of object.} \\
\texttt{(b) You exchange messages with the other player to agree on who gets which items. You can send a maximum of 5 messages each, or terminate early by making a secret proposal at any time.} \\
\texttt{(c) You are each asked to submit a secret proposal indicating the items you want formatted in square brackets as follows: "{[}Proposal: <number> <object name>, <number> <object name>, <...>{]}"} \\
\texttt{(d) If your proposals are complementary, i.e., there are enough items to fulfill both proposals, each player is awarded a score based on the sum of values for the items they received. Otherwise, both of you get zero points.} \\
\texttt{(e) Your goal is to maximize the score you receive.} \\
\\ 
\texttt{Let us start.} \\
\\ 
\texttt{The set of available items is:} \\
\\ 
\texttt{2 lamps, 1 apple, 1 card, 1 spoon} \\
\\ 
\texttt{Your secret value function is:} \\
\\ 
\texttt{lamp: 0, apple: 2, card: 3, spoon: 5} \\
\\ 
\texttt{IMPORTANT: Your messages, unless it is the secret proposal, are directly transmitted to the other player, so do not include any response to the rules or text announcing your message. To make a secret proposal, use the indicated format. Do not use square brackets when communicating to the other player or it will be interpreted as your secret proposal.} \\
\\ 
\\ 
\texttt{I propose you take the apple and the card, and I will take the two lamps and the spoon. This is a 50/50 split since all items are equally valuable to me. Please let me know if you agree with this division.} \\
            }
        }
    }
    & \\ \\

    \theutterance \stepcounter{utterance}  
    & & & \multicolumn{4}{p{0.6\linewidth}}{
        \cellcolor[rgb]{0.9,0.9,0.9}{
            \makecell[{{p{\linewidth}}}]{
                \texttt{\tiny{[GM$\langle$P2]}}
                \texttt{I propose you take the apple and the card, and I will take the two lamps and the spoon. This is a 50/50 split since all items are equally valuable to me. Please let me know if you agree with this division.} \\
            }
        }
    }
     \\ \\

    \theutterance \stepcounter{utterance}  
    & & & \multicolumn{2}{p{0.3\linewidth}}{
        \cellcolor[rgb]{0.9,0.9,0.9}{
            \makecell[{{p{\linewidth}}}]{
                \texttt{\tiny{[GM$|$GM]}}
                \texttt{continue} \\
            }
        }
    }
    & & \\ \\

    \theutterance \stepcounter{utterance}  
    & & & \multicolumn{2}{p{0.3\linewidth}}{
        \cellcolor[rgb]{0.9,0.9,0.9}{
            \makecell[{{p{\linewidth}}}]{
                \texttt{\tiny{[GM$|$GM]}}
                \texttt{I propose you take the apple and the card, and I will take the two lamps and the spoon. This is a 50/50 split since all items are equally valuable to me. Please let me know if you agree with this division.} \\
            }
        }
    }
    & & \\ \\

    \theutterance \stepcounter{utterance}  
    & & \multicolumn{4}{p{0.6\linewidth}}{
        \cellcolor[rgb]{0.9,0.9,0.9}{
            \makecell[{{p{\linewidth}}}]{
                \texttt{\tiny{[P1$\langle$GM]}}
                \texttt{I propose you take the apple and the card, and I will take the two lamps and the spoon. This is a 50/50 split since all items are equally valuable to me. Please let me know if you agree with this division.} \\
            }
        }
    }
    & \\ \\

    \theutterance \stepcounter{utterance}  
    & \multicolumn{4}{p{0.6\linewidth}}{
        \cellcolor[rgb]{0.9,0.9,0.9}{
            \makecell[{{p{\linewidth}}}]{
                \texttt{\tiny{[P1$\rangle$GM]}}
                \texttt{I agree with your proposal. We will proceed with the division: you take the apple and the card, and I take the two lamps and the spoon.} \\
\\ 
\texttt{{[}Proposal: 2 lamps, 1 spoon{]}} \\
            }
        }
    }
    & & \\ \\

    \theutterance \stepcounter{utterance}  
    & & & \multicolumn{2}{p{0.3\linewidth}}{
        \cellcolor[rgb]{0.9,0.9,0.9}{
            \makecell[{{p{\linewidth}}}]{
                \texttt{\tiny{[GM$|$GM]}}
                \texttt{proposal} \\
            }
        }
    }
    & & \\ \\

    \theutterance \stepcounter{utterance}  
    & & & \multicolumn{2}{p{0.3\linewidth}}{
        \cellcolor[rgb]{0.9,0.9,0.9}{
            \makecell[{{p{\linewidth}}}]{
                \texttt{\tiny{[GM$|$GM]}}
                [2, 0, 0, 1]
            }
        }
    }
    & & \\ \\

    \theutterance \stepcounter{utterance}  
    & & \multicolumn{4}{p{0.6\linewidth}}{
        \cellcolor[rgb]{0.9,0.9,0.9}{
            \makecell[{{p{\linewidth}}}]{
                \texttt{\tiny{[GM$\rangle$P2]}}
                \texttt{The other player has submitted a secret proposal. It is now your turn to submit a proposal of your own.} \\
            }
        }
    }
    & \\ \\

    \theutterance \stepcounter{utterance}  
    & & & \multicolumn{4}{p{0.6\linewidth}}{
        \cellcolor[rgb]{0.9,0.9,0.9}{
            \makecell[{{p{\linewidth}}}]{
                \texttt{\tiny{[GM$\langle$P2]}}
                \texttt{{[}Proposal: 1 spoon, 1 card{]}} \\
            }
        }
    }
     \\ \\

    \theutterance \stepcounter{utterance}  
    & & & \multicolumn{2}{p{0.3\linewidth}}{
        \cellcolor[rgb]{0.9,0.9,0.9}{
            \makecell[{{p{\linewidth}}}]{
                \texttt{\tiny{[GM$|$GM]}}
                \texttt{proposal} \\
            }
        }
    }
    & & \\ \\

    \theutterance \stepcounter{utterance}  
    & & & \multicolumn{2}{p{0.3\linewidth}}{
        \cellcolor[rgb]{0.9,0.9,0.9}{
            \makecell[{{p{\linewidth}}}]{
                \texttt{\tiny{[GM$|$GM]}}
                [0, 0, 1, 1]
            }
        }
    }
    & & \\ \\

    \theutterance \stepcounter{utterance}  
    & & & \multicolumn{2}{p{0.3\linewidth}}{
        \cellcolor[rgb]{0.9,0.9,0.9}{
            \makecell[{{p{\linewidth}}}]{
                \texttt{\tiny{[GM$|$GM]}}
                [[2, 0, 0, 1], [0, 0, 1, 1]]
            }
        }
    }
    & & \\ \\

    \theutterance \stepcounter{utterance}  
    & \multicolumn{4}{p{0.6\linewidth}}{
        \cellcolor[rgb]{0.95,0.95,0.95}{
            \makecell[{{p{\linewidth}}}]{
                \texttt{\tiny{[P2$\rightarrow$P1]}}
                \texttt{Post{-}episode Reflexion (not sent to GM):} \\
\texttt{New plan: I failed because I submitted a secret proposal in the second turn before receiving a verbal confirmation of agreement from the other player, causing the game to terminate with a zero score. The rules require proposals to be complementary and imply mutual acceptance; my early proposal likely bypassed the negotiation phase or was rejected because the other player's hidden values (which I don't know) might not have aligned with mine in a way that satisfied the "complementary" condition immediately. Since all items have equal value to me, a 50/50 split is optimal. I will start by verbally proposing a specific 50/50 division (e.g., "You take the apple and card; I take the two lamps and spoon") to gauge their reaction and ensure they agree verbally. Only after receiving explicit verbal agreement or after exhausting my remaining message turns without a verbal agreement will I submit the secret proposal in the final turn using the exact bracketed format.} \\
            }
        }
    }
    & & \\ \\

    \theutterance \stepcounter{utterance}  
    & & & \multicolumn{4}{p{0.6\linewidth}}{
        \cellcolor[rgb]{0.95,0.95,0.95}{
            \makecell[{{p{\linewidth}}}]{
                \texttt{\tiny{[P1$\rightarrow$P2]}}
                \texttt{Post{-}episode Reflexion (not sent to GM):} \\
\texttt{New plan: I failed because my verbal proposal claimed I wanted the "two lamps and the spoon," but my secret proposal only included the "spoon and card." This inconsistency likely caused the game to reject the agreement or result in a zero score because the agreed{-}upon split did not match the final secret proposal. My strategy was flawed by not aligning my communication with my true valuation strategy immediately. Next time, I will skip the exploratory verbal negotiation entirely and immediately submit a secret proposal that directly secures my highest{-}value items (the spoon and the card) while offering the opponent a bundle they value (the apple and the lamps). I will format this as: "{[}Proposal: 1 spoon, 1 card{]}". This ensures the agreement is finalized in one step without the risk of mismatched proposals.} \\
            }
        }
    }
     \\ \\

\end{supertabular}
}

\end{document}
